\documentclass[11pt]{article}
\usepackage[a4paper,margin=25mm]{geometry}
\usepackage{amsmath,amssymb,amsthm,booktabs}
\usepackage[colorlinks=true,urlcolor=blue,citecolor=blue]{hyperref}
\hypersetup{pdftitle={The sharp corner threshold for five-colour cubes},pdfauthor={Jason Davies}}
\newtheorem{theorem}{Theorem}
\newtheorem{lemma}[theorem]{Lemma}
\newtheorem{proposition}[theorem]{Proposition}
\theoremstyle{remark}
\newtheorem*{remark}{Remark}
\newcommand{\col}{\{0,1,2,3,4\}}
\newcommand{\mass}[2]{#1(#2)}
\title{The sharp corner threshold for five-colour cubes}
\author{Jason Davies}
\date{November 2, 2026}
\begin{document}
\maketitle
\begin{abstract}
Every collection of fifteen cubes, each displaying all five colours of a
fixed palette, contains eight that form a $2\times2\times2$ cube with
monochromatic exterior faces using those five colours. Fourteen need not
suffice: a counterexample uses just three varieties, with multiplicities
$4,5,5$. This determines the five-colour corner value listed as unknown by
Berkove, Eickemeyer and Schulman (2018). The lower bound has an elementary
proof. For the upper bound, anchored Hall deficiencies and a shared-corner
argument give 357 finite cases. Containment, budget bounds and a deletion
lemma give 279 smaller problems, each with a separately checked interval
certificate. Lean proves the interval and packed checkers sound, verifies all
279 certificates and the fourteen-cube obstruction, and connects the finite
problems to physical cube arrangements, yielding a complete formal proof of
the sharp threshold.
\end{abstract}
\noindent\small\textit{AI transparency.} The research, code and manuscript were
developed using GPT-6-Astra (OpenAI), through Codex, under the author's
direction. This working manuscript has not received independent human peer
review; Section~\ref{sec:formal} states the formal verification scope.
\normalsize

\section{Introduction}
MacMahon's coloured-cube puzzles~\cite{MacMahon} lead to a family of extremal
questions: how many face-coloured cubes guarantee that eight can be chosen
to make a larger cube with monochromatic faces? Each small cube here uses
all colours of the specified palette. Colours remain fixed, proper rotations
identify varieties, and reflections are not allowed as moves. Collections
may contain repeated varieties. Only exterior faces are constrained;
touching interior faces need not have equal colours.

Berkove, Eickemeyer and Schulman~\cite[Section 8, Table 7]{BES} tabulate the
corresponding frame numbers for palettes of two to six colours and leave
the five-colour column unknown. Their corner values for two, three, four
and six colours are respectively $11,23,11,24$. We determine the missing
five-colour value; these thresholds need not increase with the palette size.
For a $2\times2\times2$ cube the frame consists of its eight corners.
Let $N_5$ denote the least integer such that every collection of $N_5$
five-colour cubes contains a solution using all five colours on the target.
\begin{theorem}\label{thm:main}
The sharp five-colour corner threshold is $N_5=15$.
\end{theorem}
The conclusion for larger collections follows by choosing any fifteen cubes.
The theorem concerns a target that may depend on the collection; it does
not assert the ability to construct every target.

The proof has two distinct parts. A short corner argument excludes every
target for an explicit fourteen-cube collection. For the upper bound,
Hall's theorem gives linear alternatives for failure of a target.
We cover hypothetical fifteen-cube counterexamples by 357 cases. Containment
reduces these to 308, and initial budget bounds exclude another 29. A deletion
lemma reduces the totals and budgets in the remaining 279 problems, which
are excluded by finite interval certificates. The
certificates, C++ programs, geometric audits, manuscript and Lean sources
are included in the accompanying review archive.\footnote{Intended repository:
\url{https://github.com/jasondavies/macmahon-five-colour-threshold}.
Public release is pending; the review archive includes source checksums,
the base Git revision and bundled pinned project dependencies.}

A limited literature search has found no later value for this entry.
This is not a comprehensive priority claim. In particular, Rosenfeld's
October 2026 article~\cite{Rosenfeld} has not been examined in full.

\section{Geometry and Hall deficiencies}\label{sec:model}
Represent a cube by its colours on the faces
$(+x,-x,+y,-y,+z,-z)$ in a right-handed coordinate frame. A five-colour cube
has one repeated colour. There are 75 varieties: for each repeated colour,
three have its two faces opposite and twelve have them adjacent. In the
opposite case, the other four colours form two opposite pairs, giving three
possibilities. In the adjacent case, fix the two equal faces; the 24 orders
of the other colours are identified in pairs by the rotation exchanging
those equal faces. Thus there are 15 opposite-repeat varieties and 60
adjacent-repeat varieties. Each class is one orbit under colour permutations.

A corner type is the ordered triple of its incident face colours, modulo
cyclic permutation. Order the outward face normals positively, equivalently
read the triple counterclockwise from outside along the body diagonal.
This reverses the clockwise convention in~\cite{BES}, consistently on every
cube, and therefore preserves all matching and containment relations.
Repeated occurrences of a corner type are retained as distinct target
positions.

\begin{lemma}[Corner--rotation correspondence]\label{lem:corner}
A source cube fits a target corner by a proper rotation if and only if it
has the same cyclic corner type at some corner.
\end{lemma}
\begin{proof}
A proper rotation preserves the orientation of a corner frame, so it induces
a cyclic permutation of its three outward faces. Conversely, any two
positively oriented corner frames, with a prescribed cyclic alignment,
are related by a determinant-positive signed permutation of the coordinate
axes. This is a proper rotation of the cube.
\end{proof}
The palette-independent Lean version of this lemma is reused from the
pinned four-to-three-colour project~\cite{FourToThree}.

Fix a target $T$. Form a bipartite graph whose left vertices are its eight
corner positions and whose right vertices are the individual source copies.
A copy is adjacent to a position when it fits there. A matching covering
the eight positions is exactly a solution: each chosen copy can be rotated
independently, since interior faces are unrestricted. By Hall's
theorem~\cite{Hall}, failure is equivalent to a nonempty set $D$ of target
positions having fewer than $|D|$ neighbouring copies.

Index the varieties by $i=1,\ldots,75$, and let $x_i$ be their multiplicities.
For a set $S$ of varieties, write $x(S)=\sum_{i\in S}x_i$.
Let $N_T(D)$ be the varieties fitting at least one position of $D$. Then
failure of $T$ is equivalent to the existence of a cut
\begin{equation}\label{eq:hall}
  x(N_T(D))\le |D|-1,\qquad \varnothing\ne D\subseteq\{1,\ldots,8\}.
\end{equation}
There are 255 cuts per target, hence 19,125 cuts altogether.

\section{An elementary fourteen-cube obstruction}\label{sec:lower}
Consider four copies of $A$, five of $B$, and five of $C$, where
\[
 A=(0,0,1,2,3,4),\quad B=(0,1,2,3,2,4),\quad C=(0,3,1,2,4,1).
\]
Their corner types are shown below. Triples denote cyclic classes; the last
column counts positions whose type occurs in neither of the other varieties.
\begin{center}
\begin{tabular}{clc}
\toprule
Variety & Corner types & Private positions\\
\midrule
$A$ & $013,014,023,024,031,032,041,042$ & 5\\
$B$ & $022,023,034,042,122,124,132,143$ & 6\\
$C$ & $011,014,021,042,113,123,134,243$ & 6\\
\bottomrule
\end{tabular}
\end{center}
\begin{lemma}\label{lem:lower}
The collection $4A+5B+5C$ has no five-colour solution.
\end{lemma}
\begin{proof}
A solvable target must have every corner type in the displayed union.
Its only two-colour types are $022,122,011,113$. An adjacent-repeat target
therefore repeats either colour 2, with types $022,122$, or colour 1,
with types $011,113$. In either case these two corners determine the colours
on the faces at the ends of the repeated-colour edge. The two orders of the
remaining colours give a mirror pair: $B$ and its mirror, or $C$ and its
mirror. The mirror of $B$ requires $043$, and the mirror of $C$ requires
$012$; both types are absent. Only $B,C$ remain.

For an opposite-repeat target, both orientations of each occurring
unordered triple must be present. The unordered triples for which both cyclic orientations are available are
\[
  013,014,023,024,123,134.
\]
If its repeated colour is $r$, delete $r$ from the four triples containing
it. The remaining pairs must form a $K_{2,2}$ on the other four colours:
choose one face from each of their two opposite pairs. Colours 2 and 4
occur in fewer than four listed triples. For $r=1$ or $r=3$, the four pairs
have degree sequence $(3,2,2,1)$, so do not form $K_{2,2}$. For $r=0$, the
pairs are $13,14,23,24$, forcing opposite pairs $\{1,2\}$ and $\{3,4\}$.
This is $A$.

Thus only $A,B,C$ can be targets. Their private positions require respectively
at least five, six and six copies of the corresponding variety, exceeding
the available four, five and five copies. None is solvable.
\end{proof}

\section{Reduction to finite problems}\label{sec:cases}
Eight copies of any variety solve that variety, so a solution-free collection
has every multiplicity at most seven. Suppose such a collection has size 15.

\begin{lemma}[Anchored Hall cover]\label{lem:anchor}
Suppose a collection contains $m$ copies of a variety $A$ and does not solve
$A$. There is a set $K$ of its target positions such that
\begin{equation}\label{eq:anchor}
 |K|\le 7-m,\qquad
 \#\{\text{other copies fitting some position outside }K\}
       \le 7-m-|K|.
\end{equation}
Conversely, such a set certifies failure of target $A$.
\end{lemma}
\begin{proof}
Apply~\eqref{eq:hall} to $A$ and put $K=D^c$. Every copy of $A$ is a
neighbour of the nonempty set $D$. Subtract these $m$ copies from the bound
$|D|-1=7-|K|$. The converse is the same Hall inequality read backwards.
\end{proof}

First suppose the maximum multiplicity $m$ is at least two. Choose an
adjacent-repeat anchor attaining $m$ if possible; otherwise choose an
opposite-repeat anchor attaining $m$. In the latter case all adjacent-repeat
multiplicities are at most $m-1$. Relabel the colours to send the anchor to
a fixed representative of its orbit.

For each cover $K$ in Lemma~\ref{lem:anchor}, put
\[
 S(K)=N_A(K^c)\setminus\{A\},\qquad b(K)=7-m-|K|.
\]
Write $S=S(K)$ and $b=b(K)$. The case has
$x_A=m$, $x(S)\le b$, and the multiplicity bounds just described.
Additional upper bounds use two-variety solutions. If $a$ copies of $A$
and $8-a$ copies of $V$ build a target, with $a\le m$, then
$x_V\le 7-a$ in a solution-free collection. These tests are finite:
the two varieties must cover all eight target positions, $A$ must fit at
least $a$, and $V$ at least $8-a$. They are also sufficient: at most $a$
positions are unavailable to $V$,
and all of them are available to $A$. Assign these to $A$, extend the
assignment to $a$ positions that $A$ fits, and assign the remaining
$8-a$ positions to $V$.

Colour permutations and simultaneous reflection of the entire collection
preserve solvability. First identify equal pairs $(S(K),b(K))$, even when
they arise from inequivalent position subsets $K$; then retain one
representative per orbit of these pairs under the anchor stabiliser.
The reported counts are constraint-orbit counts, not position-subset counts. Transporting a cover also
transports its constrained varieties, budget and multiplicities, so no
counterexample is lost. Global reflection is used only
to transport a whole hypothetical counterexample, not as an allowed move
on an individual cube. The repeated-variety cases have the following counts,
totalling 244.
\begin{center}
\begin{tabular}{lrrrrrr}
\toprule
Anchor / maximum multiplicity & 2 & 3 & 4 & 5 & 6 & 7\\
\midrule
Adjacent-repeat & 72 & 55 & 31 & 14 & 4 & 1\\
Opposite-repeat & 26 & 21 & 11 & 6 & 2 & 1\\
Total & 98 & 76 & 42 & 20 & 6 & 2\\
\bottomrule
\end{tabular}
\end{center}

For distinct collections a different choice of anchor reduces the count.
\begin{lemma}[Shared two-colour corner]\label{lem:shared}
Among eleven adjacent-repeat cubes, two share a two-colour corner type.
If a distinct collection contains such a pair and one member is chosen as
anchor, the anchor's two two-colour positions cannot be Hall-deficient.
\end{lemma}
\begin{proof}
Each adjacent-repeat cube has exactly two such types, $aab$ and $aac$, with
$b\ne c$. There are only $5\cdot4=20$ types $aab$. Eleven cubes contribute
22 incidences, so two cubes share a type. For the two characteristic positions
of the chosen anchor, the anchor itself and the other cube are two distinct
neighbours. Their neighbourhood therefore has size at least two.
\end{proof}
If an opposite-repeat variety is present, choose it as anchor; this gives
31 constraint orbits with every multiplicity at most one. Otherwise use
Lemma~\ref{lem:shared} to choose an adjacent-repeat anchor. Omit the cover
complementary to its two characteristic positions and set all opposite-repeat
capacities to zero. Its stabiliser preserves these positions, so the omission
is compatible with the orbit reduction. It removes one of the 83
adjacent-repeat constraint orbits, leaving 82 further cases.

\begin{proposition}\label{prop:cover}
Every solution-free fifteen-cube collection is carried by colour relabelling
and, if needed, global reflection into one of the 357 exported root cases.
\end{proposition}
\begin{proof}
The preceding arguments cover repeated and distinct collections. The finite
catalogue, two-variety bounds and coverage of all anchor-cover orbits are
checked from physical rotations by the accompanying audit. It regenerates
rotations from quarter turns, separately from the signed-axis permutations
used by the case generator, and checks every exported bound and every cover
orbit. The resulting counts are
\[
  244+31+82=357.
\]
The cases need not be disjoint; completeness of the cover is what is required.
\end{proof}

The cover can be reduced without further search. If $S'\subseteq S$ and
$b'\ge b$, then $x(S)\le b$ implies $x(S')\le b'$ for nonnegative $x$.
Together with coordinatewise inclusion of the root boxes, this identifies
49 cases contained in other cases, leaving 308. Write $u$ for the root upper
bounds. In 29 retained cases, $b+u(S^c)<15$ already contradicts the prescribed
total, so only 279 cases
need further exclusion proofs. The exact map and retained case numbers are
in \path{data/q5-root-reduction.json}; Lean checks all containments and
arithmetic exclusions. The archive retains all 357 original cases.

\begin{lemma}[Hall-budget projection]\label{lem:projection}
Suppose a root has anchor multiplicity $m$, budget $x(S)\le b$, and total
$15$, with the anchor outside $S$. For $0\le s\le b$, put $G=15-b+s$
and suppose $G\ge m$. Every admissible vector $x$ has a vector $y\le x$
with the same anchor multiplicity, total $G$, budget $y(S)\le s$, and a
Hall deficiency for every target. It obeys the original upper bounds,
tightened to $\min(u_i,s)$ on $S$.
\end{lemma}
\begin{proof}
First retain only $\min(x(S),s)$ copies in $S$. This removes at most $b-s$
copies and leaves the anchor untouched, so the remaining total is at least
$G$. Delete further non-anchor copies until the total is $G$; this is possible
because $G\ge m$. Every coordinate decreases, so all upper bounds and Hall
deficiency inequalities are preserved. Nonnegativity and $y(S)\le s$ give
$y_i\le s$ for $i\in S$.
\end{proof}

For our anchored covers, $b=7-m-|K|$, so
\[
 G=m+8+|K|+s.
\]
Thus the reduced collection consists of the $m$ fixed anchor copies and
$8+|K|+s$ other copies; in particular, $G\ge m$ is automatic. Apply the lemma
after selecting and normalising the original root, retaining its bounds.
No new anchor selection is required for the smaller collection.
Of the 279 remaining roots, 277 admit an exclusion certificate with $s=0$. The other two, original
cases 68 and 69, use $s=1$ and $G=12$. Their $s=0$, $G=11$ problems actually
have solution-free collections, supplied and checked by physical matching in
the projected manifest. The projected totals range from 9 to 15. All 101 roots
whose original budget was zero retain their original certificates.

\section{Interval certificates}\label{sec:certificates}
Each case asks whether a vector of nonnegative integers $x$ can satisfy its
root bounds, prescribed total $G$, root budget and at least one Hall cut for each of
the 75 targets. We certify that no such vector exists.

Write a box as $\ell\le x\le u$. For a cut $x(S)\le r$, two necessary
conditions for feasibility in the box are
\begin{equation}\label{eq:possible}
  \ell(S)\le r,\qquad G\le r+u(S^c).
\end{equation}
The first follows from the lower bounds; the second follows from
$G=x(S)+x(S^c)$. When they hold, put
$s=r-\ell(S)$ and $g=r+u(S^c)-G$. Every satisfying vector obeys
\begin{equation}\label{eq:refine}
 x_i\le\min(u_i,\ell_i+s)\quad(i\in S),\qquad
 x_i\ge\max(\ell_i,u_i-g)\quad(i\notin S).
\end{equation}
For example, the upper bound follows by giving all other coordinates of $S$
their lower bounds; the lower bound follows by giving all other coordinates
of $S^c$ their upper bounds.

For a target, retain its cuts satisfying~\eqref{eq:possible}. If none remain,
no vector in the box can fail that target. Otherwise let $I$ and $U$ be the
intersection and union of their neighbourhoods, and let $s_*,g_*$ be the
largest corresponding $s,g$. Then all vectors in the box that fail this
target remain in the refined box
\begin{equation}\label{eq:summary}
 u'_i=\begin{cases}\min(u_i,\ell_i+s_*)&i\in I,\\u_i&i\notin I,\end{cases}
 \qquad
 \ell'_i=\begin{cases}\max(\ell_i,u_i-g_*)&i\notin U,\\\ell_i&i\in U.\end{cases}
\end{equation}
Indeed, a failing vector satisfies one of the retained cuts, and its bounds
in~\eqref{eq:refine} imply~\eqref{eq:summary}. The root budget is treated as
a family containing a single cut. The total $\sum_i x_i=G$ gives
\[
 \ell'_i=\max\!\left(\ell_i,G-\sum_{j\ne i}u_j\right),\qquad
 u'_i=\min\!\left(u_i,G-\sum_{j\ne i}\ell_j\right).
\]
A total-bound leaf closes exactly when $\sum_i\ell_i>G$ or
$\sum_i u_i<G$. A budget leaf closes when
$\ell(S)>b$ or $b+u(S^c)<G$; a target leaf closes when none of its
255 cuts satisfies~\eqref{eq:possible}. These are the checker's 77 possible
closing reasons. Refinements must tighten at least one bound and leave
nonempty coordinate intervals; a split must satisfy $\ell_i\le v<u_i$.

A certificate is a finite tree of refinements and integer splits
$x_i\le v$ or $x_i\ge v+1$. Each leaf is excluded by total bounds or by a
cut family with no surviving cut. Induction on this tree proves that its
root contains no admissible vector. The separate C++ checker reads these
instructions and verifies the deductions; it performs no search and does
not use the search program's matching implementation.

\begin{proposition}[Certificate replay]\label{prop:replay}
The separate checker accepts an exclusion certificate for each of the
357 original root cases and each of the 279 projected problems.
\end{proposition}
The two fixed archives have the following reproducible sizes.
\begin{center}
\begin{tabular}{lrr}
\toprule
 & Original & Projected\\
\midrule
Root cases & 357 & 279\\
Certificate instructions & 11,599,651 & 499,829\\
Uncompressed certificate bytes & 26,203,974 & 1,134,573\\
Compressed archive bytes & 3,516,042 & 313,094\\
\bottomrule
\end{tabular}
\end{center}
Every archived header is checked against its exported root case. Eighty
negative controls are rejected, including all 77 possible direct closing
claims against a separate $G=11$ fixture, not a $G=15$ root problem.
The witness satisfies its total, fixed anchor count, coordinate bounds and
zero root budget, and fails all 75 targets by physical matching. The fixture
is \path{data/q5-distinct-known-control.txt}; its regeneration and witness
checks are in \path{scripts/check_five_colour_distinct.py}.
An earlier matching-based exhaustive search also excluded all 357 cases,
with 772,590,466 prefix nodes; its rejecting matchings were checked as explicit
assignments of eight distinct copies. The projected archive reduces the interval instruction count by 95.69\%.
Both archives use the same checker and inference rules.

\begin{proof}[Proof of Theorem~\ref{thm:main}]
Lemma~\ref{lem:lower} gives $N_5\ge15$. A solution-free collection of fifteen
cubes would satisfy an exported root case by Proposition~\ref{prop:cover},
and would supply a Hall deficiency for every target. Containment and the initial
budget tests leave one of 279 roots. Lemma~\ref{lem:projection} supplies an
admissible vector in its projected problem, which Proposition~\ref{prop:replay}
and the sound interval deductions exclude. Hence $N_5\le15$.
\end{proof}

\section{Formalisation and reproducibility}\label{sec:formal}
The package uses Lean~4 and mathlib~\cite{Lean4,mathlib}, both at version
\texttt{v4.35.0}. Lean proves soundness of the reference interval checker,
of a faster arithmetic implementation, and of composition of bounded proof
chunks. The faster implementation encodes a multiplicity vector as
\[
 P(x)=\sum_{i=0}^{74}x_i B^i,\qquad B=1024.
\]
For a neighbourhood $S$, the weight $W(S)=\sum_{i\in S}B^{74-i}$ satisfies
\[
 x(S)=\left\lfloor\frac{P(x)W(S)}{B^{74}}\right\rfloor\bmod B.
\]
Lean proves this identity under $0\le x_i\le7$, including the absence of
carries; the coefficient bound is $75\cdot7<1024$. It also proves decoding,
coordinate updates and equivalence of the packed feasibility test to the
integer inequalities. All 19,125 supplied weights are checked against their
neighbourhood masks by the Lean kernel.

Chunk composition requires exact agreement of both boundary vectors and a
proof excluding each referenced box; the exporter is untrusted. A limit of
128 instructions gives 4,786 chunks for the projected archive, compared with
109,449 for the original archive. All 279 projected certificates have passed
kernel replay, together with their final assembly into the physical threshold
and sharpness theorems. The exact case mapping and certificate hashes are in
\path{data/q5-projected-root-manifest.json}.
The 78 declarations listed in \path{Audit.lean} and the final replay theorems
have audited dependencies confined to \texttt{propext},
\texttt{Classical.choice} and \texttt{Quot.sound}; no
\texttt{native\_decide} is used.

Lean proves completeness and uniqueness of the 75 cube varieties and checks
all 45,000 compatibility bits using the pinned corner--rotation theorem.
The physical lower-bound theorem, \path{FiveColour.Physical.lower_bound} in
\path{FiveColour/LowerBound.lean}, uses 75 deficient corner sets: 72 singletons
and three sets of sizes 5, 6 and 6. In the formal proof, this finite check replaces the
elementary target classification in Lemma~\ref{lem:lower}; catalogue completeness and rotation
transport extend it to every physical target.

Lean proves the physical Hall reduction, symmetry transport, anchor
normalisation, two-variety bounds, anchor selection and the shared-corner
argument. Finite checks certify the root capacities and cover representatives.
Containment and initial budget tests reduce the 357 roots to 279.
Lemma~\ref{lem:projection} is formalised by
\path{FiveColour.Physical.root_projection} in
\path{FiveColour/RootProjection.lean}; all concrete projection hypotheses
and coverage of the retained roots are kernel-checked.

The projected archive's integer semantics are established in
\path{FiveColour/ProjectedBridge.lean}. The identifications include all
coordinate bounds, prescribed totals, budgets and Hall-cut families.
This checks the entire admissibility predicate, not just the packed boundary
vectors. The generated module \path{FiveColour/ProjectedChunks/Threshold.lean}
assembles all 279 exclusions into the unconditional theorems
\path{FiveColour.Packed.ProjectedChunks.threshold15} and
\path{FiveColour.Packed.ProjectedChunks.sharpThreshold15}.
The original and intermediate interfaces are documented in
\path{FiveColour/README.md}. Reproduce the full proof and its axiom audit with:
\begin{verbatim}
make check WORKERS=2
make projected-replay-data WORKERS=8
make projected-replay WORKERS=8
\end{verbatim}

The review archive supplies certificates, code and reproduction commands.
A Mathlib-only comparator challenge states the puzzle geometrically.
Model equivalence is proved by \path{FiveColour.Puzzle.hasSolution_iff}
in \path{FiveColour/ComparatorBridge.lean}. The solution module applies the
completed threshold theorem and builds in Lean; external comparator checking
is in progress.
Mathematical questions left
open here include classifying all fourteen-cube obstructions and determining
the five-colour frame numbers for larger cubes.

\section*{AI disclosure and review status}
GPT-6-Astra (OpenAI), used through Codex, developed the arguments, search and
verification programs, Lean formalisation and manuscript under the human
author's direction. The author initiated and directed the research. Author
review and submission approval of this working manuscript are pending.
The inherited geometry has its own provenance in the pinned dependency.
Automated checks and AI review are not independent human peer review.

\begin{thebibliography}{9}
\raggedright
\bibitem{MacMahon} P.~A.~MacMahon,
\emph{New Mathematical Pastimes}, Cambridge University Press (1921).
\href{https://archive.org/details/newmathematicalp00macmuoft}{Digitised edition}.
\bibitem{BES} E.~Berkove, A.~Eickemeyer and M.~Schulman,
\emph{The 4-Color Cubes Puzzle}, Integers \textbf{18} (2018), A68.
\href{https://doi.org/10.5281/zenodo.10680216}{doi:\nolinkurl{10.5281/zenodo.10680216}}.
\href{https://math.colgate.edu/~integers/s68/s68.pdf}{Journal PDF}.
\bibitem{Hall} P.~Hall, \emph{On Representatives of Subsets},
Journal of the London Mathematical Society \textbf{s1-10}(1) (1935), 26--30.
\href{https://doi.org/10.1112/jlms/s1-10.37.26}{doi:\nolinkurl{10.1112/jlms/s1-10.37.26}}.
\bibitem{FourToThree} J.~Davies,
\emph{A sharp threshold for three-colour targets from four-colour cubes},
manuscript and Lean project (2026).
\href{https://github.com/jasondavies/macmahon-four-to-three/tree/6f5f04b28ac8d2156af66bbe060f2060b19b30a4}{Pinned repository revision}.
\bibitem{Rosenfeld} M.~Rosenfeld,
\emph{MacMahon Cubes and Moshe's Insanity, Puzzle 1},
Geombinatorics Quarterly \textbf{36}(2) (October 2026), starts at p.~78.
\href{https://geombina.uccs.edu/past-issues/volume-xxxvi}{Journal listing};
full text not examined.
\bibitem{Lean4} L.~de Moura and S.~Ullrich,
\emph{The Lean 4 Theorem Prover and Programming Language},
in Automated Deduction---CADE~28, LNCS \textbf{12699}, Springer (2021), 625--635.
\href{https://doi.org/10.1007/978-3-030-79876-5_37}{doi:\nolinkurl{10.1007/978-3-030-79876-5_37}}.
\bibitem{mathlib} The mathlib Community, \emph{The Lean Mathematical Library},
in Proceedings of CPP~2020, ACM (2020), 367--381.
\href{https://doi.org/10.1145/3372885.3373824}{doi:\nolinkurl{10.1145/3372885.3373824}}.
\href{https://arxiv.org/abs/1910.09336}{arXiv:1910.09336}.
\end{thebibliography}
\end{document}
